% !TEX root = main.tex

\begin{abstract}
Railway infrastructure faces risks from rail defects that can propagate into rail breaks, jeopardizing safety and reducing network efficiency. This paper develops a probabilistic forecasting framework based on a covariate-dependent Markov jump process with modified sojourn time estimates to predict defect severity progression. We test alternative generator parameterizations, each varying in structural flexibility, and benchmark them against reference models using proper scoring rules. Forecast diagnostics further assess calibration and discrimination. The Markov jump process achieves the highest prediction accuracy, with further improvements through ensemble forecasting. We demonstrate that accounting for spatial heterogeneity, such as local rail conditions, yields further improvements. Finally, we outline how predictive distributions can be embedded in a stochastic optimization program to allocate inspection and maintenance actions under budget constraints, bridging statistical modeling with practical infrastructure decision-making. The proposed framework advances probabilistic degradation modeling and provides infrastructure managers with a model framework to support predictive maintenance decision-making.
\end{abstract}
